\section{Scope, Method, and Qualifications}

\subsection*{Subject, Audience, and Companion Relationship}

This essay is a companion study to \textit{The Creature Before God}. It addresses adult readers---Catholic believers, inquirers, and philosophically serious critics---who feel the intellectual and moral weight of the ``scandal of particularity'': the question whether God's transcendence and universal salvific will render Israel's concrete election, the dated Incarnation under Pontius Pilate, instituted material sacraments, and singular personal address incoherent or provincial.

The argument proceeds in a deliberate epistemic order: first stating the four major fronts of objection at full strength, next inverting the underlying metaphysical premise by distinguishing an abstract universal from a universal cause, then unfolding the concrete biblical, Christological, sacramental, and moral economy, and finally stating the exact boundaries where natural reason yields to faith and providential mystery.

\subsection*{Authority Levels and Doctrinal Discipline}

Claims in this essay observe distinct theological registers:

\begin{itemize}
\item \textbf{Sacred Scripture} is cited as inspired revelation received within the canon and interpretive tradition of the Catholic Church.
\item \textbf{Ecumenical dogmatic definitions}---principally the Christological definition of the Council of Chalcedon (451)---state binding articles of Catholic faith.
\item \textbf{Second Vatican Council constitutions and declarations} are cited according to their specific conciliar genres: dogmatic constitutions (\textit{Dei Verbum}, \textit{Lumen gentium}), pastoral constitution (\textit{Gaudium et spes}), and declaration (\textit{Nostra aetate}).
\item \textbf{Authoritative papal and curial magisterial teaching}---including the \textit{Catechism of the Catholic Church} (promulgated by Pope St.\ John Paul II as an authoritative compendium of Catholic doctrine whose articles span defined dogmas, ordinary magisterium, and theological exposition), the papal encyclicals \textit{Redemptoris missio} and \textit{Fides et ratio}, and the Congregation for the Doctrine of the Faith's declaration \textit{Dominus Iesus}---is cited as authoritative guidance on revelation, Christ's unique mediation, sacraments, and the bounds of religious pluralism.
\item \textbf{Patristic and scholastic witnesses}---St.\ Irenaeus, St.\ Athanasius, St.\ Gregory of Nazianzus, St.\ John Chrysostom, St.\ Augustine, St.\ Leo the Great, and St.\ Thomas Aquinas---are expounded directly from identified primary editions as major Doctors of the Church, without elevating disputed school opinions (such as Thomistic versus Scotistic accounts of the metaphysical principle of individuation) into dogma.
\item \textbf{Integrative analysis} connecting these witnesses to modern historical, cosmological, and pluralistic objections represents editorial theological synthesis, carrying no magisterial or official ecclesial authority.
\end{itemize}

\subsection*{Source Boundaries, Editions, and Rights}

Every verbatim English quotation in this publication has been collated directly against locally verified, public-domain witnesses; Latin passages from St.\ Thomas Aquinas's \textit{De ente et essentia} are collated directly in the Baur--Koch Latin edition (1933) and rendered closely inline without English quotation marks:

\begin{itemize}
\item \textbf{Sacred Scripture:} English quotations follow Bishop Richard Challoner's revision of the Douay--Rheims Bible in the public-domain electronic text (Project Gutenberg eBook 1581), preserving its book names (such as \textit{Isaias} and \textit{Apocalypse}) and versification.
\item \textbf{Patristic and Scholastic Witnesses:} Verbatim English quotations of St.\ Irenaeus (\textit{Against Heresies}), St.\ Athanasius (\textit{On the Incarnation}), St.\ Gregory of Nazianzus (\textit{Oration 38}), St.\ John Chrysostom (\textit{Homilies on Matthew}), St.\ Augustine (\textit{The City of God}, trans.\ Marcus Dods, 1871; \textit{Tractates on John}), St.\ Leo the Great (\textit{Letters} and \textit{Sermons}), and St.\ Thomas Aquinas (\textit{Summa Theologiae}, Prima Pars and Tertia Pars, Fathers of the English Dominican Province) are drawn from tracked public-domain printings and CCEL/Project Gutenberg electronic witnesses, avoiding modernized web revisions where unrevised public-domain texts are available.
\item \textbf{Magisterial Sources (Paraphrase Only):} Official twentieth- and twenty-first-century Vatican texts (\textit{Dei Verbum}, \textit{Lumen gentium}, \textit{Nostra aetate}, \textit{Gaudium et spes}, \textit{Catechism of the Catholic Church}, \textit{Redemptoris missio}, \textit{Fides et ratio}, and \textit{Dominus Iesus}) remain rights-restricted under Holy See publication terms. No verbatim passages are reproduced from them; their teaching is stated in independent doctrinal paraphrase anchored to their official Latin and English paragraph numbers.
\item \textbf{Historical and Modern Interlocutors (Conceptual Engagement):} The objections associated with Celsus (via Origen's \textit{Contra Celsum}), Baruch Spinoza, Gotthold Ephraim Lessing, Immanuel Kant, S\o ren Kierkegaard, Thomas Paine, Thomas Chalmers, and John Hick---together with brief historical mentions of Boethius, Bl.\ John Duns Scotus, and Francisco Su\'arez---are engaged conceptually in analytical prose without local full-text collation or verbatim quotation.
\end{itemize}

Public-domain texts and translations remain in the public domain. Original prose, synthesis, and structure composed for this essay are offered under CC BY 4.0, carrying neither ecclesiastical approval nor critical-edition status.

\subsection*{Review State}

Primary quotations from Scripture, the Fathers, and St.\ Thomas Aquinas were directly verified against tracked local public-domain texts, audited in an independent cold review, and re-verified across all print and web build checks on 28 November 2026. No external specialist, clerical, or peer review is claimed.
